\section{Implementación práctica: cómo convertir la guía en tiempo de inferencia en un experimento reproducible}

\subsection{Los cuatro controles clave del diseño experimental}

Al avanzar desde el contenido del curso hacia su implementación en ingeniería, lo más importante no es reinventar una fórmula, sino gestionar de forma sistemática varios controles centrales.

\begin{enumerate}
    \item \textbf{Selección del sampler}: DDPM, SDE y ODE corresponden a distintos compromisos entre velocidad y estabilidad.
    \item \textbf{Programación de la intensidad de la guía}: fijar los pesos es lo más sencillo, pero suele ser menos estable que un schedule que varía según el paso temporal.
    \item \textbf{Calidad del verificador}: si el reward / verifier no es confiable, determina directamente si la dirección de la guía es correcta.
    \item \textbf{Estrategia de filtrado de candidatos}: la muestreo único, best-of-$N$ y el reordenamiento (reranking) alteran significativamente los resultados finales.
\end{enumerate}

\begin{importantbox}{Muchos experimentos fracasan no por defectos del modelo, sino por un verifier insuficiente}
Si el propio verifier no se alinea con las preferencias humanas o solo recompensa patrones locales fácilmente explotables, los modelos de difusión aprenderán a ser ``astutos''. Esto es isomórfico al problema de la explotación del reward en RLHF.
\end{importantbox}

\subsection{Plantilla recomendada para registros experimentales}

Para permitir comparaciones justas entre distintos métodos, el contenido del curso es ideal para convertirse en una plantilla unificada de registro experimental:

\begin{table}[H]
\centering
\begin{tabular}{p{0.18\textwidth} p{0.22\textwidth} p{0.42\textwidth}}
\toprule
Ítem del experimento & Campos a registrar & Descripción \\
\midrule
Modelo base & Modelo preentrenado, resolución, dominio de entrenamiento & Afecta la calidad de la estimación de Tweedie y su capacidad de transferencia \\
Configuración de muestreo & Número de pasos, sampler, semilla aleatoria & Determina la velocidad y el rango de variabilidad de los resultados \\
Configuración de guía & Tipo de reward, peso, schedule & Determina el equilibrio entre cumplimiento condicional y realismo \\
Método de verificación & Métricas automáticas, preferencias humanas, ejemplos de fallo & Evita centrarse solo en el promedio e ignorar colapsos modales \\
Costo computacional & Tiempo por muestra individual, memoria gráfica (VRAM), número de ejecuciones best-of-$N$ & Facilita la comparación de costos con métodos durante el entrenamiento \\
\bottomrule
\end{tabular}
\caption{Plantilla mínima de registro para experimentos de guía en tiempo de inferencia}
\end{table}

\begin{knowledgebox}{Por qué es importante una ``base de casos de fallo''}
Los métodos de guía de modelos de difusión suelen parecer buenos en casos promedio, pero sufren distorsiones sistemáticas ante ciertas combinaciones de condiciones o escenarios de alta resolución. Crear una base de datos específica para los ejemplos de fallo resulta más útil para la investigación posterior que simplemente reportar el valor medio.
\end{knowledgebox}

\subsection{De un solo reward a combinaciones múltiples}

En aplicaciones reales, raramente existe un único objetivo. Las imágenes deben cumplir con el texto y, además, satisfacer la consistencia geométrica, restricciones de estilo, instrucciones de edición y restricciones físicas, entre otros. Por ello, el reward total suele expresarse como:

$$r_{\text{total}}(x) = \lambda_1 r_{\text{text}}(x) + \lambda_2 r_{\text{consistency}}(x) + \lambda_3 r_{\text{physics}}(x) + \cdots$$

Esto genera dos nuevos problemas:

\begin{itemize}
    \item Las escalas de los distintos rewards difieren; sumarlas directamente permitiría que un solo ítem domine la optimización.
    \item Los objetivos de los rewards pueden ser mutuamente conflictivos; por ejemplo, una restricción geométrica más fuerte podría comprimir los detalles y la diversidad de la imagen.
\end{itemize}

\begin{warningbox}{Error más común en combinaciones múltiples de rewards: ajustar los pesos a la intuición sin fundamento}
Sin normalización del reward, recorte de gradientes o asignación por paso temporal, las distintas restricciones se interferirán entre sí, resultando en una salida ni realista ni condicionalmente válida. Una práctica más segura es que distintos rewards asuman el control principal en diferentes periodos temporales; por ejemplo, enfatizar la estructura global en etapas tempranas y los detalles y la consistencia en fases posteriores.
\end{warningbox}

\subsection{Fronteras de investigación: de la guía hand-crafted al verifier aprendido}

La tendencia implícita del taller es que, a corto plazo, la aproximación de Tweedie junto con funciones de reward manuales sigue siendo la vía más rentable; sin embargo, a largo plazo, la dirección más probable es entrenar verificadores o modelos de valor más potentes:

\begin{itemize}
    \item \textbf{Modelo de reward aprendido}: acercar el verifier a las preferencias humanas o al índice de éxito en tareas aguas abajo.
    \item \textbf{Modelo de valor temporal}: no solo evaluar el $x_0$ final, sino también el potencial futuro de los estados intermedios $x_t$.
    \item \textbf{Cálculo adaptativo}: asignar más pasos de inferencia a muestras difíciles en lugar de tratar todas las entradas por igual.
    \item \textbf{Pipeline híbrido}: durante el entrenamiento, el modelo controla la dirección general; durante la inferencia, la guía se encarga del pulido final.
\end{itemize}

\subsection{Protocolo de evaluación: no limitarse a un único indicador visual}

Aunque el curso se centra principalmente en la deducción metodológica, desde una perspectiva metodológica de investigación, la guía en tiempo de inferencia requiere especialmente evaluaciones multidimensionales. Un protocolo más completo debería incluir al menos lo siguiente:

\begin{enumerate}
    \item \textbf{Cumplimiento condicional}: error de reconstrucción del problema inverso, tasa de cumplimiento de restricciones y éxito de la tarea.
    \item \textbf{Realismo}: FID, CLIP score, puntuación subjetiva humana o métrica perceptual específica para la tarea.
    \item \textbf{Estabilidad}: si la variabilidad de los resultados es controlable bajo diferentes semillas aleatorias.
    \item \textbf{Costo}: tiempo wall-clock por muestreo individual, memoria gráfica (VRAM) y el costo adicional derivado del best-of-$N$.
\end{enumerate}

\begin{warningbox}{Reportar solo ``la mejor imagen'' exagera gravemente la eficacia del método}
La guía en tiempo de inferencia puede seleccionar fácilmente casos estéticos mediante múltiples muestreos, pero sin reportar el rendimiento promedio, la tasa de fallos y el costo computacional, no se pueden comparar de forma justa distintos métodos. Las ventajas y desventajas del escalado en tiempo de prueba deben divulgarse simultáneamente.
\end{warningbox}

\subsection{Por qué esta ruta es especialmente importante para equipos medianos y pequeños}

Un juicio realista enfatizado por el ponente es que, ante la continua elevación del umbral para entrenar modelos grandes, los métodos en tiempo de prueba ofrecen a laboratorios pequeños y medianos otra vía competitiva.

\begin{itemize}
    \item No es necesario entrenar desde cero modelos generativos a gran escala.
    \item Se pueden reutilizar directamente pesos públicos y samplers de código abierto.
    \item El foco de la investigación se desplaza hacia el diseño del verifier, la construcción del reward y la asignación de cómputo.
    \item Son más adecuados para iteraciones rápidas, ya que los ciclos de retroalimentación por experimento son mucho más cortos que los requeridos para un nuevo entrenamiento.
\end{itemize}

\begin{knowledgebox}{La comunidad de modelos de difusión podría replicar la evolución de la comunidad de LLM}
Los LLM han trazado una trayectoria muy clara: tras la comercialización de los modelos preentrenados, lo que realmente abre brechas es el post-entrenamiento, la evaluación y la orquestación en tiempo de prueba. Es probable que los modelos de difusión evolucionen siguiendo una ruta similar, y el foco competitivo futuro tenderá a inclinarse hacia el verifier, las herramientas (tooling) y el diseño de sistemas.
\end{knowledgebox}

\subsection{Resumen de la sección}

Convertir realmente la guía en tiempo de inferencia en un sistema de investigación o producto depende de unificar los registros experimentales, diseñar combinaciones estables de rewards y tratar la calidad del verifier como una prioridad absoluta. Muchos problemas que parecen pertenecer a la capa algorítmica son, en esencia, cuestiones de diseño experimental y de evaluación.
